\chapter{燃气轮机的建模分析方法}

燃气轮机是燃气-蒸汽联合循环机组中能量转换的核心设备，其主要作用是将燃料化学能经燃烧转化为高温高压燃气的热能，并进一步通过透平膨胀转化为轴功和发电功率。对于二拖一联合循环机组而言，两台燃气轮机不仅承担主要发电负荷，同时其排气也是两台余热锅炉的热源。因此，燃气轮机侧的空气流量、燃料流量、压气机出口状态、透平排气温度和排气组分等参数，会直接影响余热锅炉蒸汽产量、汽轮机进汽状态以及全厂发电供热性能。

燃气轮机主要由压气机、燃烧室、透平组成。压气机负责将环境空气压缩至燃烧所需压力，其出口压力和出口温度决定燃烧室入口状态，对压气机耗功和整机效率具有重要影响。燃烧室负责完成天然气与压缩空气的混合和燃烧，其建模重点在于燃料化学能输入、燃烧产物组成计算、燃烧室压力损失以及出口燃气状态确定。透平则将高温燃气的焓降转化为机械功，其输出功一部分用于驱动压气机，剩余部分经发电机输出，其高温排气进入余热锅炉继续释放余热，用于产生高压、中压和低压蒸汽。

在部件级建模中，压气机入口流股为大气环境下的空气，压气机出口流股为高压空气；燃烧室入口包括压缩空气和天然气燃料，出口为燃烧产物烟气；透平入口为高温烟气，出口为余热锅炉入口烟气。该划分方式与第 2 章建立的工质流股模型保持一致，能够使气体物性、质量守恒、能量平衡和㶲平衡在各部件之间连续传递。

燃气轮机基本热力过程可用布雷顿循环描述。理想布雷顿循环通常由压气机等熵压缩、燃烧室等压加热、透平等熵膨胀和排气等压放热四个过程组成。实际燃气轮机中，压气机和透平内部存在摩擦、流动分离和泄漏等不可逆损失，燃烧室存在压力损失和燃烧不完全损失，因此实际循环与理想循环之间存在明显偏差。本文在建立燃气轮机模型时，以实际测点数据为基础，将压气机、燃烧室和透平分别作为控制体进行建模，从而反映各部件内部的能量转换和㶲损失。

本文具体分析对象为某二拖一燃气-蒸汽联合循环热电联产机组中的两台燃气轮机。该机组由两台燃气轮机、两台燃气轮发电机、两套余热锅炉、一套蒸汽轮机及其发电机构成，燃气发电机组与蒸汽发电机组采用不同轴布置。燃气轮机型号为 SGT5-4000F(4)，采用单燃料天然气燃烧方式，仅用于联合循环运行。单台燃气轮机主要包括一台 15 级轴流式压气机、一个由 24 个低 NOx 燃烧器组成的环形燃烧系统、一台 4 级透平以及相关辅助系统。压气机和透平共用一根转子，并由两个支持轴承支撑，额定转速为 $3000\,\mathrm{r/min}$。

本文以全厂历史运行数据为基础构造燃气轮机模型输入。主要输入参数包括压气机进口气体状态、压气机出口气体状态、燃料气体状态、透平出口气体状态。各股气体状态包括温度、压力、流量和气体组分信息。由于燃气轮机进口空气质量流量难以测得，本文以修正后的余热锅炉烟气质量流量与燃料质量流量之差作为进口空气流量。

本文燃气轮机分析以能量分析为主。能量分析以热力学第一定律为基础，用于计算压气机耗功和等熵效率、燃料能量输入、透平输出功和等熵效率、燃气轮机发电功率和热效率等指标，评价工质能量的转化情况。此外，本文在部件建模中也定义了㶲效率等指标，为后续进一步开展㶲分析预留了接口。

\section{压气机的建模方法}

压气机是燃气轮机中完成空气压缩的主要部件，其作用是将环境空气压缩至燃烧室和透平冷却系统所需压力。压气机压缩过程需要消耗透平输出功的一部分，因此压气机耗功、压比、等熵效率和㶲效率是燃气轮机部件级分析中的重要指标。本文程序中的压气机模型以入口空气流股和实测出口空气流股为基础，并通过可调节的抽气参数近似考虑压气机冷却抽气的影响。

本文用下角标 1 表示压气机入口状态，下角标 2 表示压气机实测出口状态，下角标 $2s$ 表示与实际出口压力相同、且熵等于入口熵的等熵出口状态，下角标 $b$ 表示压气机抽气状态。

等熵出口状态 $2s$ 用于评价实际压缩过程偏离理想压缩过程的程度，满足以下条件：
\begin{equation}\label{eq:compressor-isentropic-state}
P_{2s}=P_2,
\qquad
s_{2s}=s_1.
\end{equation}

实际燃气轮机中，压气机部分空气会被抽出用于透平冷却、密封或其他辅助用途。由于运行数据中通常缺少抽气测点，本文采用参数化方式确定压气机抽气。设压气机入口空气质量流量为 $\dot{m}_1$，则抽气质量流量和主流空气质量流量分别为
\begin{equation}\label{eq:bleed-mass-flow}
\dot{m}_b = \beta_b \dot{m}_1,
\qquad
\dot{m}_{2,\mathrm{main}} = \left(1-\beta_b\right)\dot{m}_1,
\end{equation}
抽气压力和抽气比焓分别为
\begin{equation}\label{eq:bleed-state}
P_b = P_1 + \alpha_p\left(P_2-P_1\right),
\qquad
h_b = h_1 + \alpha_h\left(h_2-h_1\right),
\end{equation}
式中，$\beta_b$ 为抽气质量分数，$\alpha_p$ 为抽气压力位置系数，$\alpha_h$ 为抽气焓位置系数，三者均为可手工调节的参数。再根据 $P_b$、$h_b$ 和空气组分确定抽气温度 $T_b$、比熵 $s_b$ 和比㶲 $e_{x,b}$。

压气机压比定义为出口压力与入口压力之比 $\pi_c = P_2/P_1$。压气机等熵效率定义为等熵压缩比焓升与实际压缩比焓升之比：
\begin{equation}\label{eq:compressor-isentropic-eff}
\eta_{c,s}
=
\frac{h_{2s}-h_1}{h_2-h_1}.
\end{equation}

压气机耗功为：
\begin{equation}\label{eq:compressor-power}
\dot{W}_c
=
\dot{m}_{2,\mathrm{main}} h_2
+
\dot{m}_b h_b
-
\dot{m}_1 h_1.
\end{equation}

等熵损失功率为：
\begin{equation}\label{eq:compressor-isentropic-loss}
\dot{I}_{c,s}
=
\dot{m}_2
\left(
h_2-h_{2s}
\right).
\end{equation}

压气机㶲效率为：
\begin{equation}\label{eq:compressor-exergy-eff}
\eta_{c,ex}
=
\frac{e_{x,2s}-e_{x,1}}{e_{x,2}-e_{x,1}}.
\end{equation}

\section{燃烧室的建模方法}

燃烧室是燃气轮机中完成燃料化学能向高温燃气热能转化的部件。压气机出口空气与天然气在燃烧室内混合燃烧，燃烧产物随后进入透平膨胀做功。本文程序中的燃烧室模型以压气机出口主流空气、燃料气和燃烧产物组分为输入，通过质量守恒、能量平衡和压力恢复关系确定燃烧室出口状态。压气机抽气与高温燃气的等效冷却掺混放在透平建模部分处理。

本文用下角标 2 表示进入燃烧室的压气机出口主流空气状态，下角标 $f$ 表示燃料气状态，下角标 3 表示燃烧室出口状态。

程序首先根据天然气组分计算燃料低位发热量。对于由多种可燃组分构成的燃料气，其质量基准低位发热量可表示为
\begin{equation}\label{eq:fuel-lhv}
LHV_f
=
\frac{
\sum_i x_i M_i LHV_i
}{
M_f
},
\end{equation}
其中，$x_i$ 为燃料中第 $i$ 种组分的摩尔分数，$M_i$ 为该组分摩尔质量，$LHV_i$ 为该组分质量基准低位发热量，$M_f$ 为燃料气平均摩尔质量。

设燃烧效率为 $\eta_{\mathrm{comb}}$（默认取 $0.99$），则燃烧室内实际释放热量为
\begin{equation}\label{eq:combustion-heat-release}
\dot{Q}_{\mathrm{rel}}
=
\eta_{\mathrm{comb}} \cdot LHV_f \cdot \dot{m}_f.
\end{equation}
燃烧室压力损失通过总压恢复系数 $\sigma_{\mathrm{cc}}$（默认取 $0.95$）描述，燃烧室出口压力为 $P_3 = \sigma_{\mathrm{cc}} P_2$，对应的压力损失为 $\Delta P_{\mathrm{cc}} = (1-\sigma_{\mathrm{cc}})P_2$。当需要根据试验数据或设备资料校核燃烧室模型时，可通过调整 $\eta_{\mathrm{comb}}$ 和 $\sigma_{\mathrm{cc}}$ 改变燃烧放热量和燃烧室出口压力。

燃烧室出口烟气质量流量由空气主流质量流量和燃料质量流量相加得到：
\begin{equation}\label{eq:combustor-mass-balance}
\dot{m}_3
=
\dot{m}_{2,\mathrm{main}}
+
\dot{m}_f.
\end{equation}
忽略燃烧室与外界的散热以及动能、位能变化时，燃烧室能量平衡可写为
\begin{equation}\label{eq:combustor-energy-balance}
\dot{m}_{2,\mathrm{main}}h_2
+
\dot{m}_f h_f
+
\dot{Q}_{\mathrm{rel}}
=
\dot{m}_3 h_3.
\end{equation}
由此可得燃烧室出口比焓
\begin{equation}\label{eq:combustor-outlet-enthalpy}
h_3
=
\frac{
\dot{m}_{2,\mathrm{main}}h_2
+
\dot{m}_f h_f
+
\dot{Q}_{\mathrm{rel}}
}{
\dot{m}_{2,\mathrm{main}}+\dot{m}_f
}.
\end{equation}
已知 $P_3$、$h_3$ 和燃烧产物组分 $\boldsymbol{x}_{g}$ 后，反算燃烧室出口温度 $T_3$，并得到 $s_3$ 和 $e_{x,3}$。综上，本文燃烧室模型输出的主要参数包括燃料低位发热量 $LHV_f$、燃料热输入 $\dot{Q}_{f}=LHV_f\dot{m}_f$、实际释放热量 $\dot{Q}_{\mathrm{rel}}$、燃烧室压力损失 $\Delta P_{\mathrm{cc}}$ 以及燃烧室出口状态 $\mathcal{F}_3$。

\section{透平的建模方法}

透平是燃气轮机中将高温烟气热能转化为轴功的部件。燃烧室出口高温烟气进入透平后膨胀做功，其中一部分轴功用于驱动压气机，其余部分经发电机输出。透平模型以燃烧室出口状态、压气机抽气状态和透平实际出口状态为基础，并通过等效冷却掺混的方式将压气机抽气对透平入口状态的影响纳入考虑，从而计算透平膨胀比、实际比焓降、等熵效率、㶲效率、输出功率和等熵损失功率等指标。

本文用下角标 3 表示燃烧室出口高温烟气状态，下角标 $b$ 表示压气机抽气状态，下角标 $3c$ 表示等效冷却后的透平入口状态，下角标 4 表示透平实际出口状态，下角标 $4s$ 表示出口等熵状态。

透平冷却气的建模是透平建模的重点。冷却气在透平动叶和静叶上以连续的方式注入，形成气膜冷却。其参与透平做功，对透平出力和透平效率有较大影响，不能忽略。真实的透平冷却气建模复杂，因此常常采用等效处理方法。近似认为压气机抽气在透平入口前与燃烧室出口高温烟气完成绝热掺混，混合后的状态作为透平入口状态。该处理将透平内部连续分布的冷却掺混过程等效为入口前的一次混合，误差在可接受范围内。

% 注：透平等效冷却示意图已移除

等效掺混过程近似认为压力与燃烧室出口压力一致，即 $P_{3c}=P_3$。等效透平入口质量流量为 $\dot{m}_{3c}=\dot{m}_3+\dot{m}_b$。根据绝热混合的焓流率平衡，等效透平入口比焓为
\begin{equation}\label{eq:turbine-cooling-mix}
h_{3c}
=
\frac{
\dot{m}_3 h_3
+
\dot{m}_b h_b
}{
\dot{m}_{3c}
}.
\end{equation}
根据 $P_{3c}$、$h_{3c}$ 和燃烧产物组分 $\boldsymbol{x}_g$ 反算 $T_{3c}$，并得到 $s_{3c}$ 和 $e_{x,3c}$。该状态即为透平膨胀计算的入口状态。

透平等熵出口状态 $4s$ 用于评价实际膨胀过程偏离理想膨胀过程的程度，其定义为
\begin{equation}\label{eq:turbine-isentropic-state}
P_{4s}=P_4,
\qquad
s_{4s}=s_{3c}.
\end{equation}
程序根据 $P_{4s}$、$s_{4s}$ 和燃烧产物组分 $\boldsymbol{x}_g$ 反算 $T_{4s}$，并进一步得到 $h_{4s}$ 和 $e_{x,4s}$。

透平等熵效率定义为实际膨胀比焓降与等熵膨胀比焓降之比：
\begin{equation}\label{eq:turbine-isentropic-eff}
\eta_{t,s}
=
\frac{h_{3c}-h_4}{h_{3c}-h_{4s}}.
\end{equation}

在质量守恒条件下 $\dot{m}_4=\dot{m}_{3c}$，透平输出功率为
\begin{equation}\label{eq:turbine-power}
\dot{W}_t
=
\dot{m}_{3c}
\left(
h_{3c}-h_4
\right).
\end{equation}
等熵损失功率为
\begin{equation}\label{eq:turbine-isentropic-loss}
\dot{I}_{t,s}
=
\dot{m}_4
\left(
h_4-h_{4s}
\right).
\end{equation}
透平㶲效率采用实际输出功与理想等熵膨胀可用能降之比表示：
\begin{equation}\label{eq:turbine-exergy-eff}
\eta_{t,ex}
=
\frac{e_{x,3c}-e_{x,4}}{e_{x,3c}-e_{x,4s}}.
\end{equation}
综上，本文透平模型输出的主要参数包括等效透平入口状态 $\mathcal{F}_{3c}$、等熵出口状态 $\mathcal{F}_{4s}$、透平膨胀比 $\pi_t$、实际比焓降 $\Delta h_t$、等熵效率 $\eta_{t,s}$、㶲效率 $\eta_{t,ex}$、等熵损失功率 $\dot{I}_{t,s}$ 以及透平输出功率 $\dot{W}_t$。

\section{燃气轮机结果分析}

基于上述燃气轮机部件模型，本文对两台燃气轮机历史运行工况进行逐点计算，得到压气机、燃烧室和透平的主要热力性能参数。结果分析重点包括燃气轮机出力计算结果、燃气轮机能量分配特性以及压气机与透平效率变化规律。

\subsection{燃气轮机计算出力}

图~\ref{fig:gt1-air-flow-correction-output} 和图~\ref{fig:gt2-air-flow-correction-output} 分别给出了两台燃气轮机出力计算值与实测值的对比。当前计算方法对夏季工况的计算精度较高，但对于冬季工况偏差较大。这可能是由于冬季工况点较少，修正系数的拟合参数主要来源于夏季工况，导致冬季工况的修正系数预测不准确。此外，冬季环境温度较低，燃气轮机性能对空气流量变化更敏感，因此同样的空气流量误差可能导致更大的功率残差。未来可以通过增加冬季工况数据点、引入环境温度作为修正参数或采用非线性回归方法来进一步提升修正模型的适用性和准确性。对比两台燃机，燃机1的计算误差更小，这可能是由于其运行工况较为稳定，且空气流量修正函数较为适配。

\begin{figure}[!htb]
  \centering
  \includegraphics[width=0.8\textwidth]{chapter3/1燃机1出力.png}
  \caption{\#1 燃气轮机出力计算值与实测值对比}
  \label{fig:gt1-air-flow-correction-output}
\end{figure}

\begin{figure}[!htb]
  \centering
  \includegraphics[width=0.8\textwidth]{chapter3/1燃机2出力.png}
  \caption{\#2 燃气轮机出力计算值与实测值对比}
  \label{fig:gt2-air-flow-correction-output}
\end{figure}

\FloatBarrier

\subsection{压气机和透平效率分析}

图~\ref{fig:chapter3-gt1-compressor} 和图~\ref{fig:chapter3-gt2-compressor} 对比了两台燃气轮机压气机的性能计算结果。压气机效率主要反映实际压缩过程相对于等熵压缩过程的接近程度。随着燃气轮机出力升高，压气机通常运行在更接近设计工况的区域，压比和空气流量增加，流动损失相对减小，因而压气机等熵效率整体趋于稳定并保持在较高水平。冬季工况下压气机效率较低，这可能是冬季空气密度较大，体积流量较小，使得压气机偏离设计工况导致的。随燃机出力增大，空气流量相应增大，压气机效率相应升高，这与压气机性能曲线的特性一致。

透平效率反映高温烟气膨胀做功过程的有效程度。透平入口温度、膨胀比、冷却空气掺混比例和排气压力都会影响透平效率。低温环境下机组通常能够在较高负荷或更接近设计匹配的区域运行，透平膨胀比、流量系数及气动匹配状态有所改善，因此计算得到的透平效率相对较高。需要指出的是，实际运行数据中的透平效率由温度、压力、燃料量和烟气流量等测点反算得到，冬季烟气流量修正、比热取值及测点误差也可能对效率结果产生影响。

\begin{figure}[!htb]
  \centering
  \includegraphics[width=0.8\textwidth]{chapter3/3燃机1压气机.png}
  \caption{\#1 燃气轮机压气机效率}
  \label{fig:chapter3-gt1-compressor}
\end{figure}

\begin{figure}[!htb]
  \centering
  \includegraphics[width=0.8\textwidth]{chapter3/4燃机2压气机.png}
  \caption{\#2 燃气轮机压气机效率}
  \label{fig:chapter3-gt2-compressor}
\end{figure}

\begin{figure}[!htb]
  \centering
  \includegraphics[width=0.8\textwidth]{chapter3/3燃机1透平.png}
  \caption{\#1 燃气轮机透平效率}
  \label{fig:chapter3-gt1-turbine}
\end{figure}

\begin{figure}[!htb]
  \centering
  \includegraphics[width=0.8\textwidth]{chapter3/4燃机2透平.png}
  \caption{\#2 燃气轮机透平效率}
  \label{fig:chapter3-gt2-turbine}
\end{figure}

\FloatBarrier

\subsection{燃气轮机能量分配和占比分析}

图~\ref{fig:chapter3-gt1-energy-detail} 和图~\ref{fig:chapter3-gt2-energy-detail} 分别给出了两台燃气轮机的能量分配详情。燃料低位热输入进入燃烧室后，一部分经透平膨胀转化为燃气轮机轴功并由发电机输出，一部分以排气显热形式进入余热锅炉，剩余部分表现为压气机耗功、燃烧室压力损失、透平不可逆损失以及其他散热和机械损失。对于联合循环机组而言，燃气轮机排气能量并非无效损失，而是蒸汽循环的热源。

图~\ref{fig:chapter3-gt1-energy-share} 和图~\ref{fig:chapter3-gt2-energy-share} 进一步给出了燃气轮机能量累计占比随实际出力的变化。发电功率占比约为 $0.28$--$0.36$，且随燃机出力升高略有增加，说明高负荷工况下燃料输入向净输出功的转化比例有所提高；压气机耗功占比约为 $0.38$--$0.43$，是燃气轮机内部功率消耗的主要部分；尾气能量及其他损失占比约为 $0.22$--$0.32$，随出力升高总体略有下降，但仍保留了较高的余热利用潜力。高出力区间部分散点离散程度增大，可能与空气流量测量误差与修正偏差、排气温度波动和测点误差有关。

\begin{figure}[!htb]
  \centering
  \includegraphics[width=0.8\textwidth]{chapter3/2GT1能量分配.png}
  \caption{\#1 燃气轮机能量分配详情}
  \label{fig:chapter3-gt1-energy-detail}
\end{figure}

\begin{figure}[!htb]
  \centering
  \includegraphics[width=0.8\textwidth]{chapter3/2GT2能量分配.png}
  \caption{\#2 燃气轮机能量分配详情}
  \label{fig:chapter3-gt2-energy-detail}
\end{figure}

\begin{figure}[!htb]
  \centering
  \includegraphics[width=0.8\textwidth]{chapter3/2-1GT1能量占比.png}
  \caption{\#1 燃气轮机能量占比}
  \label{fig:chapter3-gt1-energy-share}
\end{figure}

\begin{figure}[!htb]
  \centering
  \includegraphics[width=0.8\textwidth]{chapter3/2-1GT2能量占比.png}
  \caption{\#2 燃气轮机能量占比}
  \label{fig:chapter3-gt2-energy-share}
\end{figure}

\FloatBarrier

\section{本章小结}

本章建立了燃气轮机部件级性能分析模型，分别对压气机、燃烧室和透平进行能量分析，并提出基于功率残差的空气流量修正方法。在此基础上，对两台燃气轮机的历史运行工况进行逐点计算，得到空气流量修正效果、部件效率和能量分配等结果。分析结果表明燃气轮机功率计算有较高的精度；压气机和透平效率反映了其在不同负荷下的工况特征；能量分配揭示了燃料输入在发电输出和排气余热之间的分配关系。上述分析为后续余热锅炉、汽轮机以及联合循环整机性能评价提供了燃气轮机侧部件级基础。
